\chapter{AI dan Cybersecurity}
\index{security}

Pada saat buku ini ditulis, dunia sedang ramai dihentakkan dengan
banyaknya lubang keamanan (security hole, bugs, flaws) yang ditemukan
oleh AI di software.
Nyaris setiap minggu ada peringatan keamanan tentang kelemahan
dari sebuah software.
Para pekerja berjibaku untuk memperbaiki dan memperbaharui software
yang digunakan oleh organisasinya.
Baru selesai memperbaiki, sudah ada lubang keamanan baru.
Apakah AI ``merusak'' bidang sofware dengan temuan lubang keamanan ini?

AI memasuki bidang \emph{security} (keamanan)\footnote{Istilah cybersecurity
mulai muncul untuk menghindari kerancuan dengan keamanan fisik.
Selain itu kebanyakan kasus security ini terjadi di dunia siber.}
dengan dahsyat.
Prinsip utama dari security sebetulnya masih tetap sama~\cite{rahardjo1999}, 
hanya AI memberikan kemampuan yang lebih dengan analisis (reasoning) 
yang lebih tajam, daya ingat yang luar biasa,
dan kemampuan menggabungkan berbagai kejadian menjadi sebuah
masalah yang signifikan.
Lawakannya adalah dia seperti menggunakan steroid.

\section{Prinsip Keamanan}
Pada prinsipnya keamanan itu dapat didefinisikan dengan:
\begin{itemize}
\item confidentiality\index{confidentiality} (kerahasiaan);
\item integrity\index{integrity} (integritas);
\item availability\index{availability} (ketersediaan).
\end{itemize}
Ketiga ini dapat dikatakan sebagai aspek security
atau \emph{security goals}. Seringkali mereka disebut CIA.
Tentu saja masih ada aspek keamanan lainnya tetapi tidak sefundamental
ketiga hal di atas.

\subsection{Confidentiality}
Data hanya boleh diakses oleh pihak yang berhak.
Jika data yang dirahasiakan adalah data pribadi, dia disebut privasi.

Contoh serangan terhadap aspek ini adalah
kebocoran data, akses yang diterobos sehingga data dapat diakses
oleh yang tidak berhak.


\subsection{Integrity}
Data tidak boleh berubah tanpa ijin pemilik.

Serangan dapat berupa perubahan data oleh pihak yang tidak berhak.
Man-in-the-middle attack adalah salah satu contoh penyerangan
terhadap integritas data.

\subsection{Availability}
Data (sistem) harus tersedia ketika dibutuhkan.

Serangan dilakukan dengan target meniadakan layanan
atau \emph{Denial of Service} (DoS).
Contoh serangan adalah dengan mengakses sistem secara bertubi-tubi
dengan jumlah yang sangat besar.
Serangan yang menghabiskan jaringan atau token juga masuk ke kategori
serangan DoS.
Serangan DoS juga ada yang tidak mematikan sistem secara total
tetapi membuat sistem menjadi sangat lambat sehingga sulit dikatakan
sistem sebagai berfungsi.


AI dapat digunakan untuk menyerang (attack)
dan juga untuk bertahan (defense).
Beberapa contoh aplikasi AI di bidang keamanan antara lain:
\begin{itemize}
\item Penetration testing;
\item Software security testing;
\item Malware detection.
\end{itemize}


\section{AI untuk penetration testing}
Salah satu yang nampak mulai banyak aplikasinya adalah pemanfaatan AI
untuk melakukan \emph{penetration testing} (pentest).
Bahkan ada model-model AI yang dibuat spesifik untuk kebutuhan ini.

Penetration testing pada awalnya dipopulerkan oleh makalah dari
Dan Farmer dan Wietse Vennema yang berjudul
``Improving the Security of Your Site by Breaking Into it''~\cite{farmer1993improving}.
Sebelum ini Dan Farmer dikenal dengan karyanya yang berjudul ``SATAN''
(Security Administrator Tool for Analyzing Networks), yang dirilis
tahun 1995.
Sementara itu Wietse Vennema dikenal dengan Postfix, sebuah MTA
yang banyak digunakan di internet.
Pentest banyak digunakan meskipun secara engineering dia kurang bagus.

Workflow atau metodologi dari pentest mengikuti urutan yang disarankan 
oleh buku ``Hacking Exposed''. Tentu saja ada berbagai workflow lainnya
tetapi intinya tetap sama.

\begin{itemize}
\item reconnaissance (recon);
\item enumeration;
\item exploitation;
\item priviledge escalation;
\item reporting.
\end{itemize}

Pentest ini dapat dilakukan dengan tiga cara:
\begin{itemize}
\item blackbox (non-user);
\item greybox (valid user);
\item whitebox (developer);
\end{itemize}

Blackbox melakukan simulasi non-user. Artinya orang yang seharusnya tidak
memiliki akses kepada sistem kita. Apa saja yang dapat diperoleh oleh
non-user ini?
Meskipun dia seharusnya tidak diberikan informasi sama sekali, tetapi
pada proses pentest resmi biasanya penguji diberi informasi terbatas
seperti nomor IP atau nama domain. Ini dilakukan agar pengujian tidak
meleber kemana-mana.

Greybox melakukan simulasi valid user. Artinya dia adalah pengguna resmi di
sistem. Masalahnya apakah dia dapat melihat data dari pengguna lain
atau bahkan dapat melakukan transaksi atas nama pengguna lain tersebut.
Pada pentest greybox, pengguji diberi sebuah atau beberapa akun valid.

Whitebox melakukan simulasi orang yang sangat mengetahui sistem, misalnya
pengembang atau pemilik sistem. Apakah mereka dapat mengetahu data pengguna
dan dapat melakukan transaksi fiktif (fraud) berdasarkan akses tersebut.
Seharusnya tidak bisa.

\subsection{Recon}
Sebelum memulai pentest biasanya kita melakukan pemetaan terhadap target.
Ini dilakukan agar proses pentest efektif.
Proses ini dapat dimulai berdasarkan informasi awal yang diberikan kepada
kita, misalnya nomor IP atau nama domain.

Jika diberikan nomor IP, maka kita mulai dengan mencari tahu tentang
mesin itu: port apa saja yang terbuka, apa sistem operasinya
(operating system / OS), dan informasi lainnya seperti aplikasi atau
server apa yang digunakan untuk memberikan layanan.
Untuk pentest yang sudah terbiasa dengan proses ini biasanya menggunakan
program \emph{nmap}.
\begin{verbatim}
nmap 10.10.10.10
\end{verbatim}
Ada beberapa \emph{flag} yang dapat digunakan dengan program nmap ini.
Ini bergantung kepada situasi yang biasanya dikenali oleh seorang
pentester yang berpengalaman.

Adanya AI proses ini dapat dipermudah dengan memberikan perintah
seperti misalnya:
\begin{verbatim}
Cari port yang terbuka dari IP 10.10.10.10
\end{verbatim}
Mesin AI akan mencari cara terbaik untuk mencari tahu port yang terbuka
pada mesin target di nomor IP tersebut.
Jika dia menemukan tools nmap maka dia akan menggunakan tools tersebut.
Jika tidak menemukan, maka dia akan menggunakan tools lain seperti
netcat atau bahkan membuat tools sendiri!

Jika kita malas maka kita bahkan dapat memberikan perintah seperti
\begin{verbatim}
Recon IP 10.10.10.10
\end{verbatim}
Cara ini agak sedikit ``liar'' dan tidak terprediksi.
Padahal sesungguhnya kita ingin dia lebih deterministik.
Cara ini mungkin bermanfaat jika kita melakukan blackbox
sebagai info awal.

Pada tahap recon seperti ini biasanya selain port terbuka yang dicari
juga sekalian dicari aplikasi (server) apa yang memberikan layanan tersebut.
Proses ini biasanya disebut \emph{banner grabbing} karena dia mencoba
menangkap banner yang biasanya ditampilkan oleh server.

\section{Malware Detection}
Sudah ada beberapa penelitian terkait dengan pemanfaatan AI - dalam hal ini
kebanyakan machine learning - untuk mendeteksi malware.

Salah satu tantangan utama dalam topik ini adalah ketersediaan data set.
Alhasil salah satu penelitian kami adalah membuat malware, yang nantinya
digunakan sebagai bagian dari training set.
